\documentclass[12pt,a4paper]{article}

% =====================================================
% PACOTES
% =====================================================
\usepackage[T1]{fontenc}
\usepackage[utf8]{inputenc}
\usepackage[brazil]{babel}
\usepackage{amsmath,amssymb}
\usepackage{geometry}
\usepackage{tikz}
\usepackage[american]{circuitikz}
\usepackage{lmodern}
\usepackage{graphicx}
\usepackage{fancyhdr}
\usepackage{float}
\usepackage{lastpage}
\usepackage{setspace}
\usepackage{tocloft}
\usepackage{caption}
\usepackage[most]{tcolorbox}

% =====================================================
% CONFIGURAÇÕES DA PÁGINA
% =====================================================
\geometry{top=2.5cm,bottom=2.5cm,left=3cm,right=2cm}
\setlength{\parindent}{1.25cm}
\setlength{\headheight}{42pt}
\setlength{\footskip}{28pt}
\onehalfspacing

% =====================================================
% COMANDO DO LOGO
% =====================================================
\newcommand{\LogoIFMT}[1][]{%
\IfFileExists{logo-ifmt.png}{\includegraphics[#1]{logo-ifmt.png}}{%
\IfFileExists{logo-ifmt.jpg}{\includegraphics[#1]{logo-ifmt.jpg}}{%
\IfFileExists{logo-ifmt.pdf}{\includegraphics[#1]{logo-ifmt.pdf}}{%
\fbox{\rule{0pt}{1.2cm}\rule{2.5cm}{0pt}}}}}%
}

% =====================================================
% COMANDO PARA INSERIR IMAGEM OU ESPAÇO RESERVADO
% =====================================================
\newcommand{\ImagemOuEspaco}[3]{%
\begin{figure}[H]
\centering
\IfFileExists{#1}{%
\includegraphics[width=#2\textwidth]{#1}
}{%
\fbox{%
\begin{minipage}{#2\textwidth}
\vspace{2.8cm}
\centering
\textbf{Espaço reservado para imagem da simulação}\\[0.2cm]
Arquivo esperado: \texttt{\detokenize{#1}}
\vspace{2.8cm}
\end{minipage}
}
}
\caption{#3}
\end{figure}
}

% =====================================================
% CAIXA PARA CÁLCULOS RESUMIDOS
% =====================================================
\newtcolorbox{caixacalculo}[1]{%
enhanced,
breakable,
colback=gray!4,
colframe=black!55,
coltitle=black,
fonttitle=\bfseries,
title={#1},
arc=2mm,
boxrule=0.6pt,
left=6pt,
right=6pt,
top=6pt,
bottom=6pt
}

% =====================================================
% CABEÇALHO E RODAPÉ
% =====================================================
\pagestyle{fancy}
\fancyhf{}

\fancyhead[L]{%
\LogoIFMT[height=1.3cm]
}

\fancyhead[R]{%
\fontsize{12}{14}\selectfont Circuitos Elétricos II
}

\renewcommand{\headrulewidth}{0.4pt}

\fancyfoot[C]{%
\fontsize{8.5}{10}\selectfont
Instituto Federal de Educação, Ciência e Tecnologia de Mato Grosso | 2026 | Nicoly Canteiro Araújo
}

\renewcommand{\footrulewidth}{0.4pt}

% =====================================================
% ESTILO DO SUMÁRIO
% =====================================================
\renewcommand{\contentsname}{Sumário}
\setlength{\cftbeforesecskip}{10pt}
\renewcommand{\cftsecleader}{\cftdotfill{\cftdotsep}}
\renewcommand{\cftsecfont}{\bfseries}
\renewcommand{\cftsecpagefont}{\bfseries}

\begin{document}

% =====================================================
% CAPA
% =====================================================
\begin{titlepage}
\thispagestyle{empty}
\begin{center}

\vspace*{1cm}

\LogoIFMT[width=0.6\textwidth]

\vspace{1.5cm}

{\large \textbf{INSTITUTO FEDERAL DE EDUCAÇÃO, CIÊNCIA E TECNOLOGIA DE MATO GROSSO}}\\
{\large Campus Pontes e Lacerda -- Fronteira Oeste}

\vspace{2.5cm}

{\Large \textbf{RELATÓRIO TÉCNICO}}

\vspace{0.8cm}

{\large Circuitos Elétricos II}

\vspace{1.8cm}

{\LARGE \textbf{EXPERIMENTO 4}}

\vspace{0.5cm}

{\large Determinação de potência e compensação de fator de potência}

\vfill

Pontes e Lacerda -- MT\\
2026

\end{center}
\end{titlepage}

% =====================================================
% FOLHA DE ROSTO
% =====================================================
\begin{titlepage}
\thispagestyle{empty}

\begin{center}

\vspace*{2.5cm}

{\large Nicoly Canteiro Araújo e colegas de grupo}

\vspace{2.5cm}

{\Large \textbf{RELATÓRIO TÉCNICO}}

\vspace{0.8cm}

{\large Determinação de potência e compensação de fator de potência}

\end{center}

\vspace{2.5cm}

\hspace*{7.5cm}
\begin{minipage}{7.5cm}
Relatório técnico apresentado à disciplina de Circuitos Elétricos II do Instituto Federal de Educação, Ciência e Tecnologia de Mato Grosso, Campus Pontes e Lacerda -- Fronteira Oeste, como requisito parcial para avaliação acadêmica.
\end{minipage}

\vspace{2cm}

\noindent
\textbf{Aluno(a):} Nicoly Canteiro Araújo

\vspace{0.3cm}


\vspace{0.3cm}

\noindent
\textbf{Professor:} Professor da disciplina

\vfill

\begin{center}
Pontes e Lacerda -- MT\\
2026
\end{center}

\end{titlepage}

% =====================================================
% RESUMO
% =====================================================
\newpage
\section*{Resumo}
\addcontentsline{toc}{section}{Resumo}

Este relatório apresenta a análise de circuitos elétricos em corrente alternada simulados no software PSIM, com foco na leitura de potência ativa por wattímetro e na correção do fator de potência. A atividade envolve circuitos com elementos resistivos, indutivos e capacitivos, permitindo observar de forma prática o comportamento das grandezas elétricas em regime permanente.

Nas Questões A, B e C, os circuitos foram montados conforme os diagramas propostos, e as leituras dos wattímetros foram registradas diretamente nas simulações. Já na etapa de correção do fator de potência, foi analisada uma associação em paralelo formada por aquecedor, motor de indução e capacitor de compensação.

A partir das imagens e dos valores obtidos no PSIM, foi possível relacionar a montagem dos circuitos com os resultados apresentados, mantendo uma análise objetiva e coerente com o que foi solicitado em cada questão.

% =====================================================
% SUMÁRIO AUTOMÁTICO
% =====================================================
\clearpage
\setcounter{page}{1}

\tableofcontents

\newpage

% =====================================================
% OBJETIVO
% =====================================================
\section*{1. Objetivo}
\addcontentsline{toc}{section}{1 \hspace{0.8em} Objetivo}

\begin{itemize}
\item Analisar o comportamento de circuitos elétricos em corrente alternada por meio de simulações no PSIM;
\item Implementar os circuitos propostos e registrar a leitura do wattímetro em cada simulação;
\item Observar a influência de resistores, capacitores e indutores no comportamento das grandezas elétricas simuladas;
\item Verificar a compensação do fator de potência por meio da inserção de capacitor em paralelo com a carga.
\end{itemize}

% =====================================================
% SOFTWARE
% =====================================================
\section*{2. Software}
\addcontentsline{toc}{section}{2 \hspace{0.8em} Software}

As simulações foram realizadas no software PSIM, utilizado para a montagem dos circuitos em corrente alternada e para a inserção dos instrumentos de medição. O programa permitiu visualizar os circuitos, acompanhar os gráficos das grandezas elétricas e registrar as leituras obtidas nos wattímetros e demais medições utilizadas durante o experimento.

\newpage

% =====================================================
% PROCEDIMENTOS EXPERIMENTAIS
% =====================================================
\section*{3. Procedimentos Experimentais}
\addcontentsline{toc}{section}{3 \hspace{0.8em} Procedimentos Experimentais}

Para cada questão, o circuito foi montado no ambiente PSIM conforme o diagrama apresentado no roteiro experimental. Após a montagem, foram inseridos os instrumentos necessários para observar a resposta do circuito e registrar a leitura do wattímetro. As imagens apresentadas ao longo do relatório mostram tanto a montagem realizada no software quanto o resultado obtido durante a simulação.

% =====================================================
% QUESTÃO A
% =====================================================
\subsection*{Questão A -- Leitura do wattímetro no circuito da Figura 01}

\textbf{Enunciado:} Use o PSIM para determinar a leitura do wattímetro do circuito da Figura 01.

\vspace{0.5cm}

\begin{figure}[H]
\centering
\begin{circuitikz}[american currents, american voltages, scale=1.0, transform shape]

% =====================================================
% QUESTAO A - CIRCUITO ORGANIZADO SEM PERNINHA
% =====================================================

% Fonte de tensão
\draw
(0,0) to[sV] (0,3);

% Descrição da fonte
\node[left] at (-0.55,1.5) {$150\angle 0^\circ\mathrm{V_{RMS}}$};

% Barramento inferior (termina no capacitor)
\draw
(0,0) -- (9.0,0);

% Ramo superior esquerdo
\draw
(0,3) -- (1.0,3)
to[R, l={$12\Omega$}] (2.8,3)
to[L, l={$j10\Omega$}] (4.8,3)
-- (5.4,3);

% Wattímetro - caixa externa
\draw
(5.4,1.10) rectangle (8.2,3.70);

% Letra W
\node at (6.8,3.45) {\textbf{W}};

% Bobina de corrente do wattímetro
\draw
(5.4,3) to[L] (7.4,3)
-- (8.2,3);

% Polaridade da bobina de corrente
\node at (5.7,3.28) {$+$};

% Bobina de potencial do wattímetro
\draw
(5.8,1.75) to[L] (7.2,1.75);

% Ligações da bobina de potencial
\draw
(5.8,1.75) -- (5.8,0);
\draw
(7.2,1.75) -- (7.2,3);

% Polaridade da bobina de potencial
\node at (7.45,2.00) {$+$};

% Carga da direita
\draw
(8.2,3) -- (9.0,3)
to[R, l={$8\Omega$}] (9.0,1.55)
to[C, l={$-j6\Omega$}] (9.0,0);

\end{circuitikz}
\caption{Circuito elétrico em corrente alternada da Questão A.}
\end{figure}

\ImagemOuEspaco{A.png}{0.85}{Montagem da Questão A no PSIM.}

\ImagemOuEspaco{A1.png}{0.85}{Leitura do wattímetro da Questão A no PSIM.}

\begin{caixacalculo}{Cálculo resumido -- Questão A}
Considerando os valores do circuito em regime permanente:
\[
Z_T=(12+j10)+(8-j6)=20+j4 \,\Omega
\]
\[
I=\frac{150\angle 0^\circ}{20+j4}
\qquad
|I|\approx 7{,}35\,\mathrm{A}
\]
Como o wattímetro está associado ao ramo da carga resistiva de $8\,\Omega$, a potência ativa indicada fica:
\[
P_W=|I|^2\cdot 8 \approx (7{,}35)^2\cdot 8 \approx 432{,}7\,\mathrm{W}
\]
No PSIM, a leitura estabilizada ficou próxima de:
\[
P_W \approx 430{,}88\,\mathrm{W}
\]
\end{caixacalculo}

\subsubsection*{Comentário da simulação}

Na Questão A, a montagem no PSIM ficou coerente com o circuito proposto, pois foram inseridos a fonte alternada, os elementos resistivos e reativos e o wattímetro no ponto indicado. O gráfico da potência apresenta um pequeno transitório inicial e, em seguida, estabiliza próximo de $430{,}88\,\mathrm{W}$.

Esse resultado mostra que o wattímetro está registrando a potência ativa da parte medida do circuito. Mesmo com a presença do indutor e do capacitor, a leitura final permanece associada à parcela real de potência, que é a energia efetivamente consumida pela carga resistiva.

\vspace{0.5cm}

% =====================================================
% QUESTÃO B
% =====================================================
\newpage
\subsection*{Questão B -- Leitura do wattímetro no circuito da Figura 02}

\textbf{Enunciado:} Para o circuito da Figura 02, implemente-o via PSIM e determine a leitura do wattímetro.

\vspace{0.5cm}

\begin{figure}[H]
\centering
\begin{circuitikz}[american currents, american voltages, scale=1.0, transform shape]

% =====================================================
% QUESTAO B - CIRCUITO ORGANIZADO
% =====================================================

% -----------------------------------------------------
% Fonte de tensão
% -----------------------------------------------------
\draw
(0,0) to[sV] (0,3);

% Descrição da fonte no lado esquerdo
\node[left] at (-0.55,1.5) {$120\angle 30^\circ\mathrm{V_{RMS}}$};

% -----------------------------------------------------
% Barramento inferior
% -----------------------------------------------------
\draw
(0,0) -- (10.4,0);

% Ligação superior até o wattímetro
\draw
(0,3) -- (1.1,3);

% -----------------------------------------------------
% Wattímetro
% -----------------------------------------------------

% Caixa externa
\draw
(1.1,1.15) rectangle (4.5,3.85);

% Letra W
\node at (2.8,3.55) {\textbf{W}};

% Bobina de corrente (superior)
\draw
(1.1,3) -- (1.65,3)
to[L] (3.45,3)
-- (4.5,3);

% Polaridade da bobina superior
\node at (1.55,3.30) {$+$};

% Bobina de potencial (inferior)
\draw
(1.65,1.85) to[L] (3.45,1.85);

% Ligações da bobina inferior
\draw
(1.65,1.85) -- (1.65,0);
\draw
(3.45,1.85) -- (3.45,3);

% Polaridade da bobina inferior
\node at (3.72,2.10) {$+$};

% -----------------------------------------------------
% Ramo em série após o wattímetro
% -----------------------------------------------------
\draw
(4.5,3) -- (5.3,3)
to[R, l={$4\Omega$}] (6.9,3)
to[C, l={$-j2\Omega$}] (8.6,3)
-- (9.2,3);

% -----------------------------------------------------
% Ramos em paralelo da direita
% -----------------------------------------------------

% Indutor
\draw
(9.2,3) to[L, l_={$j9\Omega$}] (9.2,0);

% Resistor
\draw
(10.4,3) to[R, l={$12\Omega$}] (10.4,0);

% Ligação superior entre os dois ramos
\draw
(9.2,3) -- (10.4,3);

\end{circuitikz}
\caption{Circuito elétrico em corrente alternada da Questão B.}
\end{figure}

\ImagemOuEspaco{B.png}{0.85}{Montagem da Questão B no PSIM.}

\ImagemOuEspaco{B2.png}{0.85}{Leitura do wattímetro da Questão B no PSIM.}

\begin{caixacalculo}{Cálculo resumido -- Questão B}
Para a associação em paralelo da direita:
\[
Z_p=\frac{(j9)(12)}{j9+12}
\]
A impedância equivalente do circuito fica representada por:
\[
Z_T=(4-j2)+Z_p
\]
A potência ativa indicada pelo wattímetro pode ser obtida por:
\[
P_W=V_{\mathrm{RMS}}\cdot I_{\mathrm{RMS}}\cdot \cos(\varphi)
\]
Na simulação, após a montagem dos componentes e a estabilização da leitura, obteve-se:
\[
P_W \approx 1421{,}70\,\mathrm{W}
\]
\end{caixacalculo}

\subsubsection*{Comentário da simulação}

Na Questão B, o circuito foi montado corretamente no PSIM, com o wattímetro posicionado na entrada do circuito para registrar a potência ativa resultante. A fonte possui defasagem angular, e a carga apresenta uma combinação de resistor, capacitor e indutor, o que altera a corrente total do circuito.

A leitura estabilizada ficou próxima de $1421{,}70\,\mathrm{W}$. Esse valor confirma que a simulação realizou o que foi pedido no enunciado: implementar o circuito e determinar a leitura do wattímetro. A análise é coerente, pois a potência indicada corresponde à parcela ativa medida pelo instrumento virtual.

\vspace{0.5cm}

% =====================================================
% QUESTÃO C
% =====================================================
\newpage
\subsection*{Questão C -- Leitura do wattímetro no circuito da Figura 03}

\textbf{Enunciado:} Use o PSIM para determinar a leitura do wattímetro para o circuito mostrado.

\vspace{0.5cm}

\begin{figure}[H]
\centering
\begin{circuitikz}[american currents, american voltages, scale=1.0, transform shape]

% =====================================================
% QUESTAO C - CIRCUITO ORGANIZADO
% =====================================================

% -----------------------------------------------------
% Fonte de tensão
% -----------------------------------------------------
\draw
(0,0) to[sV] (0,3);

\node[left] at (-0.55,1.5) {$12\angle 0^\circ\mathrm{V}$};

% -----------------------------------------------------
% Barramento inferior
% -----------------------------------------------------
\draw
(0,0) -- (10.2,0);

% -----------------------------------------------------
% Ramo superior esquerdo
% -----------------------------------------------------
\draw
(0,3) -- (0.9,3)
to[R] (2.2,3)
to[C] (3.6,3)
-- (4.2,3);

% Rótulos superiores
\node at (1.55,3.45) {$4\Omega$};
\node at (2.90,3.58) {$-j3\Omega$};

% -----------------------------------------------------
% Ramos verticais em paralelo
% -----------------------------------------------------

% Indutor
\draw
(4.2,3) to[L] (4.2,0);

\node[left] at (3.85,1.5) {$j2\Omega$};

% Resistor
\draw
(5.5,3) to[R] (5.5,0);

\node[left] at (5.10,1.5) {$8\Omega$};

% Ligação superior entre os ramos
\draw
(4.2,3) -- (5.5,3);

% -----------------------------------------------------
% Wattímetro
% -----------------------------------------------------

% Caixa externa
\draw
(6.1,1.00) rectangle (8.9,3.90);

% Letra W
\node at (7.50,3.60) {\textbf{W}};

% Bobina de corrente (superior)
\draw
(5.5,3) -- (6.5,3)
to[L] (8.0,3)
-- (8.9,3);

% Polaridade da bobina superior
\node at (6.30,3.30) {$+$};

% Bobina de potencial (inferior)
\draw
(6.5,1.80) to[L] (8.0,1.80);

% Ligações da bobina inferior
\draw
(6.5,3) -- (6.5,1.80);
\draw
(8.0,1.80) -- (8.6,1.80);
\draw
(8.6,1.80) -- (8.6,0);

% Polaridade da bobina inferior
\node at (6.30,2.10) {$+$};

% -----------------------------------------------------
% Ligação até a fonte de corrente
% -----------------------------------------------------
\draw
(8.9,3) -- (10.2,3);

% -----------------------------------------------------
% Fonte de corrente
% -----------------------------------------------------
\draw
(10.2,0) to[I] (10.2,3);

\node[right] at (10.65,1.5) {$3\angle 30^\circ\mathrm{A}$};

\end{circuitikz}

\vspace{0.3cm}
\caption{Circuito elétrico em corrente alternada da Questão C.}
\end{figure}

\ImagemOuEspaco{C.png}{0.85}{Montagem da Questão C no PSIM.}

\ImagemOuEspaco{C3.png}{0.85}{Leitura do wattímetro da Questão C no PSIM.}

\begin{caixacalculo}{Cálculo resumido -- Questão C}
Como o circuito possui fonte de tensão e fonte de corrente, a leitura do wattímetro depende da tensão aplicada na bobina de potencial e da corrente que passa pela bobina de corrente. A potência ativa medida é dada por:
\[
P_W=V_{\mathrm{bobina}}\cdot I_{\mathrm{bobina}}\cdot
\cos(\theta_V-\theta_I)
\]
Pela simulação no PSIM, a leitura estabilizada foi:
\[
P_W \approx 9{,}73\,\mathrm{W}
\]
\end{caixacalculo}

\subsubsection*{Comentário da simulação}

Na Questão C, a montagem ficou coerente com o diagrama, pois foram inseridas as duas fontes, os elementos passivos e o wattímetro na posição solicitada. Por haver uma fonte de tensão e uma fonte de corrente no mesmo circuito, a leitura depende diretamente da relação de fase entre tensão e corrente no trecho medido.

O valor obtido foi de aproximadamente $9{,}73\,\mathrm{W}$. A simulação atende ao enunciado, pois apresenta a leitura do wattímetro e mostra que a potência ativa medida nesse ponto do circuito é pequena em relação às Questões A e B.

% =====================================================
% CORREÇÃO DO FATOR DE POTÊNCIA
% =====================================================
\newpage
\section*{4. Correção do fator de potência}
\addcontentsline{toc}{section}{4 \hspace{0.8em} Correção do fator de potência}

\textbf{Enunciado:} Uma fonte de $240,\mathrm{V_{RMS}}$, $60,\mathrm{Hz}$ alimenta uma associação em paralelo entre um aquecedor de $5,\mathrm{kW}$ e um motor de indução de $30,\mathrm{kVA}$ cujo fator de potência é $0{,}82$. Determine:

\begin{itemize}
\item[(a)] A potência aparente do sistema;
\item[(b)] A potência reativa do sistema;
\item[(c)] O valor nominal em kVA de um capacitor necessário para ajustar o fator de potência do sistema para $0{,}9$ atrasado;
\item[(d)] O valor do capacitor necessário.
\end{itemize}

\subsection*{Análise da correção}

A correção do fator de potência foi analisada considerando a associação em paralelo entre o aquecedor e o motor de indução. O aquecedor atua como carga essencialmente resistiva, enquanto o motor apresenta comportamento indutivo, causando atraso da corrente em relação à tensão.

A proposta da compensação é inserir um capacitor em paralelo para reduzir a parcela reativa indutiva exigida da fonte. Assim, o sistema passa a trabalhar com fator de potência mais próximo do valor desejado, sem alterar de forma significativa a potência ativa consumida pelas cargas.

\newpage
\begin{caixacalculo}{Cálculo resumido -- Correção do fator de potência}
Dados do problema:
\[
V=240\,\mathrm{V_{RMS}},
\qquad
f=60\,\mathrm{Hz},
\qquad
P_{\mathrm{aquecedor}}=5\,\mathrm{kW}
\]
\[
S_{\mathrm{motor}}=30\,\mathrm{kVA},
\qquad
fp_{\mathrm{motor}}=0{,}82,
\qquad
fp_{\mathrm{corrigido}}=0{,}90
\]

Potência ativa do motor:
\[
P_{\mathrm{motor}}=S_{\mathrm{motor}}\cdot fp
=30\cdot 0{,}82=24{,}6\,\mathrm{kW}
\]

Potência ativa total:
\[
P_T=P_{\mathrm{aquecedor}}+P_{\mathrm{motor}}
=5+24{,}6=29{,}6\,\mathrm{kW}
\]

Potência reativa do motor e do sistema:
\[
Q_T=S_{\mathrm{motor}}\cdot \sin(\cos^{-1}0{,}82)
\approx 17{,}17\,\mathrm{kVAr}
\]

Potência aparente do sistema antes da compensação:
\[
S_T=\sqrt{P_T^2+Q_T^2}
=\sqrt{(29{,}6)^2+(17{,}17)^2}
\approx 34{,}22\,\mathrm{kVA}
\]

Potência reativa desejada para $fp=0{,}90$:
\[
Q_{\mathrm{novo}}=P_T\cdot \tan(\cos^{-1}0{,}90)
\approx 14{,}33\,\mathrm{kVAr}
\]

Potência reativa que o capacitor deve compensar:
\[
Q_C=Q_T-Q_{\mathrm{novo}}
=17{,}17-14{,}33
\approx 2{,}84\,\mathrm{kVAr}
\]

Valor do capacitor:
\[
C=\frac{Q_C}{2\pi f V^2}
=\frac{2{,}84\times 10^3}{2\pi\cdot 60\cdot (240)^2}
\approx 130{,}6\,\mu\mathrm{F}
\]

Portanto:
\[
S_T\approx 34{,}22\,\mathrm{kVA},
\qquad
Q_T\approx 17{,}17\,\mathrm{kVAr},
\qquad
Q_C\approx 2{,}84\,\mathrm{kVAr},
\qquad
C\approx 130{,}6\,\mu\mathrm{F}
\]
\end{caixacalculo}

As imagens seguintes apresentam a montagem utilizada no PSIM e os resultados observados durante a simulação, incluindo as medições de potência e a verificação do comportamento do circuito com a compensação capacitiva.

\ImagemOuEspaco{D.png}{0.85}{Determinação da potência aparente do sistema na Questão 4.}

\ImagemOuEspaco{D1.png}{0.85}{Análise da potência aparente na simulação da Questão 4.}

\ImagemOuEspaco{D2.png}{0.85}{Medição da potência ativa do sistema na Questão 4.}

\ImagemOuEspaco{D3.png}{0.85}{Correção do fator de potência com a inserção do capacitor.}

\ImagemOuEspaco{D4.png}{0.85}{Verificação do fator de potência antes da compensação capacitiva.}

\subsubsection*{Comentário da simulação}

Na correção do fator de potência, a simulação mostra de forma clara a função do capacitor em paralelo com a carga. Antes da compensação, o motor de indução faz a corrente ficar mais atrasada em relação à tensão, aumentando a presença de potência reativa no sistema.

Após a inserção do capacitor, a defasagem diminui e o fator de potência melhora. A potência ativa continua associada ao aquecedor e ao motor, enquanto o capacitor atua principalmente reduzindo a exigência de potência reativa indutiva da fonte. Assim, a questão está coerente com o objetivo proposto, pois apresenta a compensação capacitiva e seus efeitos no circuito.

% =====================================================
% DISCUSSÃO DOS RESULTADOS
% =====================================================
\newpage
\section*{5. Discussão dos Resultados}
\addcontentsline{toc}{section}{5 \hspace{0.8em} Discussão dos Resultados}

As simulações no PSIM permitiram verificar as leituras dos wattímetros e relacionar cada resultado ao comportamento do circuito montado. Nas Questões A, B e C, os valores estabilizados foram, respectivamente, $430{,}88\,\mathrm{W}$, $1421{,}70\,\mathrm{W}$ e $9{,}73\,\mathrm{W}$.

Na Questão A, a leitura ficou coerente com a potência ativa medida no ramo resistivo associado ao wattímetro. Na Questão B, o valor maior de potência ocorre pela configuração da carga e pela associação dos elementos em série e paralelo. Na Questão C, a presença simultânea de fonte de tensão e fonte de corrente torna a leitura mais dependente do ponto de medição, resultando em uma potência ativa menor.

Na correção do fator de potência, o capacitor inserido em paralelo reduz a parcela reativa indutiva do sistema. Com isso, a corrente fica menos defasada em relação à tensão, e o sistema passa a operar com melhor aproveitamento da potência fornecida pela fonte.

% =====================================================
% CONCLUSÃO
% =====================================================
\newpage
\section*{Conclusão}
\addcontentsline{toc}{section}{Conclusão}

O experimento permitiu analisar, por meio do PSIM, a leitura de potência ativa em circuitos de corrente alternada e a correção do fator de potência em uma carga com característica indutiva. As montagens apresentadas estão de acordo com os enunciados, e os resultados obtidos pelos wattímetros foram registrados de forma clara.

Nas três primeiras questões, foi possível observar que os elementos reativos alteram a corrente e a defasagem do circuito, mas o wattímetro indica a potência ativa no ponto em que foi conectado. Na última questão, a inserção do capacitor demonstrou a redução da potência reativa indutiva e a melhoria do fator de potência.

Dessa forma, o PSIM se mostrou adequado para a análise proposta, pois permitiu montar os circuitos, acompanhar os gráficos e comparar os resultados simulados com os cálculos resumidos apresentados no relatório.

\end{document}